% ---
% id: statistics/binomial-box/prop:rule-bernstein
% kind: prop
% title: "The Bernstein endpoint rule"
% summary: >-
%   Bernstein's inequality with upward-rounded rational bounds on the
%   logarithm and the root gives closed-form endpoints, an admissible rule,
%   monotone in the allocation when every rounding is. Its unrounded
%   half-width falls strictly in n as n^(−1/2), and it is never narrower than
%   the exact rule.
% prominence: standard
% status: implemented
% uses: [def:endpoint-rule, def:interval, lem:bernstein-tail, lem:log-enclosure, lem:guard, statistics/def:conservative-evaluation]
% mentions: [prop:rule-exact, thm:monotonicity-limits, req:rule-in-force, req:honest-model-record]
% assumes: []
% refs: []
% provenance: newly-derived
% implements: certiqs_key.statistics.rules:BernsteinRule
% lean:
% notes: "The rule is defined inside the proposition, as the exact rule is inside def:interval: one definition of what a rule is (def:endpoint-rule) and one claim per rule. The conservative-evaluation hypothesis is for the rounding direction of the evaluation."
% ---

The emitted interval of the
\tref[exact binomial quantile rule]{exact rule} is not monotone in $n$
(Theorem~\ref{thm:monotonicity-limits}\,(v)). The Bernstein rule is the alternative whose
unrounded half-width decreases strictly with $n$; that its integer-rounded
boxes are nested in $n$ is not claimed. Its endpoints come in closed form
from Bernstein's inequality, with no tail evaluated, and every
logarithm and root is bounded by rationals rounded in the direction that can
only widen the interval, the direction a \tref{conservative evaluation}
(Definition~\ref{def:conservative-evaluation}) admits.
A request selects it (Requirement~\ref{req:rule-in-force}), and $j$ and $s$
are recorded with the rule (Requirement~\ref{req:honest-model-record}).

\begin{proposition}\label{prop:rule-bernstein}
Assume a \tref{conservative evaluation}
(Definition~\ref{def:conservative-evaluation}). Fix an integer $j\ge1$ and a
precision $s\ge0$ (parameters of the evaluation). The \emph{Bernstein rule}
$R_{\mathrm{B}}$ assigns to
$(n,p,\varepsilon^{-},\varepsilon^{+})$, with $v\coloneqq np(1-p)$:
\begin{itemize}[leftmargin=1.9em]
\item if $p=0$: $R^{-}_{\mathrm{B}}=R^{+}_{\mathrm{B}}=0$; if $p=1$:
$R^{-}_{\mathrm{B}}=R^{+}_{\mathrm{B}}=n$;
\item otherwise, for each sign separately: if $\varepsilon^{-}=0$ then
$R^{-}_{\mathrm{B}}=0$, and if $\varepsilon^{+}=0$ then
$R^{+}_{\mathrm{B}}=n$; if $\varepsilon^{\mp}>0$, let $L^{\mp}$ be the
rational upper bound $m(u-1)$ of \eqref{eq:log-rational-bounds} on
$\ln(1/\varepsilon^{\mp})$ with $m=2^{j}$ and $u$ the upper root enclosure
of \eqref{eq:isqrt-bounds} at precision $s$, let
\begin{equation}\label{eq:bernstein-endpoint-t}
  t^{\mp}\ \coloneqq\ \text{a rational }\ \ge\
  \frac{L^{\mp}}{3}+\sqrt{\Big(\frac{L^{\mp}}{3}\Big)^{2}+2vL^{\mp}}
\end{equation}
obtained by rounding the square root upward with \eqref{eq:isqrt-bounds},
and put
\begin{equation}\label{eq:bernstein-endpoints}
  R^{-}_{\mathrm{B}}\ \coloneqq\ \max\big\{0,\ \lfloor np-t^{-}\rfloor\big\},
  \qquad
  R^{+}_{\mathrm{B}}\ \coloneqq\ \min\big\{n,\ \lceil np+t^{+}\rceil\big\}.
\end{equation}
\end{itemize}
Then:
\begin{enumerate}[label=(\roman*),leftmargin=2.2em]
\item $R_{\mathrm{B}}$ is an \tref{admissible endpoint rule} in the sense of
Definition~\ref{def:endpoint-rule}\,(a)--(c), computed from $n$, $p$ and the
allocation by finitely many operations on integers and rationals (integer
square roots included) and no others; Lemma~\ref{lem:guard} therefore
applies to it;
\item $R_{\mathrm{B}}$ is \tref{monotone in the allocation}, provided every
rounding in $L^{\mp}$ and $t^{\mp}$ is a non-decreasing function of its
real argument. The integer-square-root rounding of \eqref{eq:isqrt-bounds}
is \emph{not} such a function of a rational argument as written (its excess over the
root depends on the denominator), but it becomes one when the argument is
first rounded up to the fixed dyadic grid $2^{-s}\mathbb{Z}$, i.e.\ replaced
by $\lceil y\,2^{s}\rceil/2^{s}$ before the root is taken; the evaluation
takes $u$ from \eqref{eq:isqrt-bounds} applied to that grid point (with
$d=2^{s}$ fixed) at every step of Lemma~\ref{lem:log-enclosure}'s
iteration, rounds the square root in \eqref{eq:bernstein-endpoint-t} the
same way, and takes $L^{\mp}$ and $t^{\mp}$ themselves as the least
multiples of $2^{-s}$ at or above the values so obtained. Every such step
is a composition of non-decreasing maps, and every rounding is upward, so
(i) is unaffected;
\item \emph{(half-width)} for $0<p<1$ and $\varepsilon^{\mp}>0$, the
\tref{guarded endpoints} of $R_{\mathrm{B}}$ satisfy
\begin{equation}\label{eq:bernstein-width}
  0\ \le\ q^{+}(\hat c)-p\ \le\ \frac{t^{+}}{n}+\frac1n,
  \qquad
  0\ \le\ p-q^{-}(\hat c)\ \le\ \frac{t^{-}}{n}+\frac1n,
\end{equation}
where, for the unrounded value of \eqref{eq:bernstein-endpoint-t},
\begin{equation}\label{eq:bernstein-halfwidth}
  \frac{t^{\mp}}{n}
  \ =\ \frac{L^{\mp}}{3n}+\sqrt{\Big(\frac{L^{\mp}}{3n}\Big)^{2}
        +\frac{2p(1-p)L^{\mp}}{n}}
\end{equation}
is strictly decreasing in $n$ and is $O(n^{-1/2})$ at fixed $p$ and
$\varepsilon^{\mp}$. The rule never emits an interval narrower than the raw
quantile interval of the exact rule of Definition~\ref{def:interval}
(Proposition~\ref{prop:rule-exact}\,(iii)).
\end{enumerate}
\end{proposition}

\begin{proof}
\emph{(i).} Definition~\ref{def:endpoint-rule}\,(c) and (b) hold by the
first two clauses of the rule. For (a) at $p\in\{0,1\}$: the count is $0$
respectively $n$ almost surely, so with $R^{\mp}_{\mathrm{B}}=0$
respectively $n$ both tails of \eqref{eq:rule-admissible} vanish and the
values lie in $\{0,\dots,n\}$. For (a) at $0<p<1$: the values
\eqref{eq:bernstein-endpoints} lie in $\{0,\dots,n\}$ by the clipping. If
$\varepsilon^{+}=0$ then $R^{+}_{\mathrm{B}}=n$ and
$\Pr[\Ncount{}>n]=0\le\varepsilon^{+}$. If $\varepsilon^{+}>0$: by
Lemma~\ref{lem:log-enclosure}, $L^{+}\ge\ln(1/\varepsilon^{+})$, so
$e^{-L^{+}}\le\varepsilon^{+}$; by Lemma~\ref{lem:bernstein-tail} with
$L=L^{+}$, every $t\ge t(L^{+})$ has $\Pr[\Ncount{}-np\ge t]\le e^{-L^{+}}$,
and $t^{+}\ge t(L^{+})$ by \eqref{eq:bernstein-endpoint-t}. If
$R^{+}_{\mathrm{B}}=n$ the tail is $0$. Otherwise
$R^{+}_{\mathrm{B}}=\lceil np+t^{+}\rceil\ge np+t^{+}$, so the event
$\{\Ncount{}>R^{+}_{\mathrm{B}}\}$ is contained in
$\{\Ncount{}-np>t^{+}\}\subseteq\{\Ncount{}-np\ge t^{+}\}$ and
$\Pr[\Ncount{}>R^{+}_{\mathrm{B}}]\le e^{-L^{+}}\le\varepsilon^{+}$. The
lower side is the mirror image: if $R^{-}_{\mathrm{B}}=0$ the tail
$\Pr[\Ncount{}<0]$ is $0$; otherwise
$R^{-}_{\mathrm{B}}=\lfloor np-t^{-}\rfloor\le np-t^{-}$, so
$\{\Ncount{}<R^{-}_{\mathrm{B}}\}\subseteq\{\Ncount{}-np<-t^{-}\}
\subseteq\{\Ncount{}-np\le-t^{-}\}$, whose probability is at most
$e^{-L^{-}}\le\varepsilon^{-}$ by the second inequality of
\eqref{eq:bernstein}. The computation uses the rational arithmetic of the
inputs, the integer square roots of \eqref{eq:isqrt-bounds} for $u$ and for
the root in \eqref{eq:bernstein-endpoint-t} (applied to the rational
$(L^{\mp}/3)^{2}+2vL^{\mp}$), and floor, ceiling, maximum and minimum.

\emph{(ii).} Fix all arguments but $\varepsilon^{+}$ and let
$\varepsilon^{+}_1\le\varepsilon^{+}_2$, both positive. Then
$1/\varepsilon^{+}_1\ge1/\varepsilon^{+}_2$. With the argument first
rounded up to the grid, $y\mapsto\lceil y\,2^{s}\rceil$ is non-decreasing,
and for a \emph{fixed} denominator $d=2^{s}$ the bound
$u=(\operatorname{isqrt}(c\,d\,4^{s})+1)/(d\,2^{s})$ of
\eqref{eq:isqrt-bounds} is non-decreasing in the integer numerator $c$
because $\operatorname{isqrt}$ is non-decreasing on the integers; so each
step of the root iteration is non-decreasing in its input, and $u$ is
non-decreasing in $1/\varepsilon^{+}$. The final upward rounding to the
grid is non-decreasing too, so $L^{+}_1\ge L^{+}_2$; the map
$L\mapsto L/3+\sqrt{(L/3)^{2}+2vL}$ is increasing in $L\ge0$ and the upward
rounding of the root is monotone, so $t^{+}_1\ge t^{+}_2$; hence
$R^{+}_{\mathrm{B},1}\le R^{+}_{\mathrm{B},2}$ by
\eqref{eq:bernstein-endpoints}, which is the required non-increase of
$\varepsilon^{+}\mapsto R^{+}_{\mathrm{B}}$. At $\varepsilon^{+}_1=0$ the
value $n$ is the largest possible, so the monotonicity extends to the zero
budget. The lower side is the mirror image ($R^{-}_{\mathrm{B}}$
non-decreasing in $\varepsilon^{-}$, with the value $0$ at
$\varepsilon^{-}=0$ the smallest possible).

\emph{(iii).} The left inequalities of \eqref{eq:bernstein-width} are
\eqref{eq:ordering}. For the right ones: by \eqref{eq:guard} and
\eqref{eq:bernstein-endpoints}, $k^{+}(\hat c)=\max\{R^{+}_{\mathrm{B}},
\lceil np\rceil\}\le\lceil np+t^{+}\rceil\le np+t^{+}+1$ (using
$\lceil np\rceil\le\lceil np+t^{+}\rceil$ and $n\le\lceil np+t^{+}\rceil$
when the clip at $n$ is active), so $q^{+}(\hat c)-p\le(t^{+}+1)/n$;
symmetrically $k^{-}(\hat c)=\min\{R^{-}_{\mathrm{B}},\lfloor np\rfloor\}
\ge\lfloor np-t^{-}\rfloor\ge np-t^{-}-1$ when the clip at $0$ is inactive,
and $k^{-}(\hat c)\ge0>np-t^{-}-1$ when it is active (then
$\lfloor np-t^{-}\rfloor\le0$, so $np-t^{-}<1$); hence
$p-q^{-}(\hat c)\le(t^{-}+1)/n$. Dividing \eqref{eq:bernstein-endpoint-t} by
$n$ gives \eqref{eq:bernstein-halfwidth}; each of its two terms is strictly
decreasing in $n$ (the first as $1/n$, the second as the square root of a sum
of a term in $1/n^{2}$ and a term in $1/n$, both decreasing, at fixed
$p\in(0,1)$ and $L^{\mp}>0$), and it is $O(n^{-1/2})$ because the first term
is $O(n^{-1})$ and the radicand is $O(n^{-1})$. The comparison with the exact
rule: the exact rule's $\kappa^{+}$ is the minimum of the set in
\eqref{eq:kappaplus} (Definition~\ref{def:interval}), and
$R^{+}_{\mathrm{B}}\in\{0,\dots,n\}$ lies in that set by (i), so
$R^{+}_{\mathrm{B}}\ge\kappa^{+}$; likewise $\kappa^{-}$ is the maximum of the
set in \eqref{eq:kappaminus}, in which $R^{-}_{\mathrm{B}}$ lies by (i), so
$R^{-}_{\mathrm{B}}\le\kappa^{-}$; the guard preserves both inequalities.
\end{proof}
